%----------------------------------------------------------------
%
%  File    :  project-webtech.tex
%
%  Author  :  Martin Rabensteiner
% 
%  Created :  31 Mar 2026
% 
%  Changed :  03 Jul 2026
% 
%----------------------------------------------------------------

\chapter{Modern Web Technologies}
\label{chap:webtech}

Web technologies are a steadily changing field to keep track with new
technical possibilities and optical or interactivity affordances. This
chapter covers some topics related to the development of ListMerger.

The main browser vendors united in 2004 as Web Hypertext
Application Technology Working Group (WHATWG), but the definition of the
HTML standard was in charge of the World Wide Web Consortium (W3C) for
some decades. In 2018, the browser vendors became mutinous and only one
year later, both signed an agreement where W3C is giving the vendors
more say \parencite{cimpanu-whatwg-w3c}. This led to more
standardization among different browsers, and a faster drive of
innovation, like the improvements mentioned as follows.




\section{Living HTML}

The Hyper Text Markup Language (HTML) standard was first developed and
published in the early 1990's by Tim Berners-Lee at CERN and later
handed over to the new founded World Wide Web Consortium
(W3C)~\parencite{whatwg-livingstandard}. The last major version (HTML4)
was published in 1998, and after some intermediate versions, the work on
HTML5 started by W3C in 2004. HTML5 was never published in a finished
version. Since 2019, W3C and WHATWG are maintaining the standard
together as a living standard, and publish updates continuously in the
publishing browers.




\section{Modern CSS}

Along with the HTML modernisation, the Cascading Style Sheets (CSS)
became a living standard as well \parencite{w3c-css}. CSS is used to
define to optical style of a website. Modern solutions include for
example containers queries \parencite{mdn-cssquieries} to differentiate
between different screen sizes or use styles as selectors, or colour
mixing functions \parencite{mdn-colorfunctions}. With the evolving
functionality of CSS queries, more and more use cases can be implemented
either in CSS or JavaScript, or with selectors interdependent with both
\parencite{css-states-js}.




\section{Modern JavaScript}

JavaScript is essential in modern web development for interactivity and
essential for frontend and moren and more backend development
\parencite{shukla-jsfullstack}. It enables client-site computations and
modifying the Document Object Model (DOM), which describes the structure
and content of a web page. It was introduced in 1995, soon available in
all usual browsers, and is and still today an important component of
them. As dynamic web pages become more and more popular, it was even
more necessary to manipulate the content directly on the client.
The current standard of JavaScript is ECMAScript 2026, the 17th edition
of ECMA-262 \parencite{ecma2026}.

% TS
TypeScript is a script language written by Microsoft
\parencite{typescript}. It is based on JavaScript and extends it with
type signature and compile-time type checking to better support the
development of large applications. For production, TypeScript code is
transpiled to JavaScript. Some concepts like classes were already taken
over by the ECMA standard and are now available in plain JavaScript.

% Node.js
With Node.js, first published in 2009, it became more and more popular
to use JavaScript also as a server-side programming language
\parencite{nodejs}. JavaScript can be therefore used in the backend as
well as in the frontend, creating new synergies.
%




\section{The Fading of jQuery}

jQuery is a still widley-used frontend framework and reached its peak of
usage on around 78 \% of all websites tracked in 2022, but is steadily
decreasing since then, with circa 70 \% in March 2026
\parencite{w3techs-js-statistics}. By time, the outstanding functions of
jQuery got more and more redundant. Several use cases and
functionalities were integrated directly in HTML, CSS, or, mostly,
JavaScript. For example, the accordion can now be replaced with the
\elname{details} element (available since 2020)
\parencite{mdn-html-details}. JavaScript introduced selector functions
like \vname{querySelector()}, \vname{querySelectorAll()} and some other
for HTML element in 2015 \parencite{mdn-js-queryselector}. This made the
classic jQuery DOM selection with the {\vname{\$()}} function obsolete.




\section{JavaScript Package Managers}

The node package manager (NPM) was introduced in 2009, to make the
handling with node modules in the development process and their
publishment among developers more convenient. It is an open source
project and now in the hands of GitHub \parencite{npm-website}. Package
variables and dependencies are defined in the \fname{package.json} file.

PNPM (performant NPM) introduces a more performant approach for package
managing than NPM \cite{pnpm-website}. Instead of downloading and
storing dependencies for every project on its own, they are managed
globally on the system, and symbolic links are created from the project
to the main storage. This saves download time and storage space.
Further, PNPM is capable of managing different versions of Node. For
other package managers, a separate Node Version Manager is needed.
Disadvantages arise when multiple disks are used on a system, as the
symlinks used by PNPM are not capable of that.

To measure the popularity of package managers, \textcite{maj-article}
discusses different ways to count the usage. Different ways of download
counting have the problem of different sources, and packager managers
with more frequent updates will have higher download counts. The
observation of the most popular GitHub repositories is considered a
meaningful measurement. According to the last iteration
\cite{maj-repository} and compared with the last three years, NPM is
steadily used in about two third of projects. The usage of Yarn classic
has decreased from 29 \% to 18 \% over the same time frame. PNPM has a
stagnant share of about 6 \%.

Another currently uprising package manager is bun, which claims to
install dependencies a few times faster than the tools mentioned before
\parencite{bun-pm}. Its toolkit provides alternatives for Node.js, Vite,
and other JavaScript development tools as well.




\section{Mustache Templates}
\label{sec:mustache}

Mustache is a logic-less template format \parencite{mustache}.
Characteristic are the double \vname{{{}}} brackets as placeholders. For
example, \vname{{{variable}}} is replaced with the value of
\vname{variable}. For-each loops and simple if conditions can be mocked
with Mustache, even though there is no logic in the template. Mustache
libraries are available for around 50 programming languages, including
JavaScript, Python, PHP, or C++. It is popular to structure small or
large parts of content with its template style.
